
\Problem{Urgencias Hospitalarias}

\Statement

Se quiere simular el funcionamiento de unas urgencias
hospitalarias. Los pacientes llegan a urgencias con una cierta
prioridad (1, 2 o 3, donde 1 es la más alta) y un identificador.
Cuando el médico está libre, se atiende al paciente con la
prioridad más alta (el número más pequeño).

\medskip

En urgencias hay una sola consulta, de modo que solo un paciente
puede estar dentro a la vez. La entrada consiste en una secuencia
de eventos de dos tipos:

\begin{itemize}

\item \verb|arriba <id> <prioridad> <timestamp>|: llega un
paciente con identificador \verb|id|, prioridad \verb|prioridad|
(1, 2 o 3) y en el instante \verb|timestamp|.

\item \verb|seguent <timestamp>|: el paciente que está en la
consulta sale, y el paciente en la cola con más prioridad
entra a la consulta en el instante \verb|timestamp|.

\end{itemize}

\medskip

Cada vez que un paciente entra a la consulta (evento
\verb|seguent|), se escribe su identificador. Cuando se lee el
evento \verb|seguent|, el paciente que entra es el que tiene la
prioridad más alta entre los que están esperando. Si hay empate,
se atiende al que llegó antes.

\medskip

Al final de la entrada, se escribe el tiempo de espera medio de
los pacientes atendidos para cada prioridad (1, 2 y 3). El tiempo
de espera de un paciente es la diferencia entre el momento en que
entra a la consulta y el momento en que llegó.

\Input

Una secuencia de eventos, uno por línea. Los eventos tienen el
formato descrito más arriba. La entrada garantiza que los eventos
tienen sentido: nunca hay más de un paciente en la consulta, y
cuando se hace \verb|seguent| hay al menos un paciente esperando.

\Output

Para cada evento \verb|seguent|, se escribe el identificador del
paciente que entra a la consulta.

\medskip

Al final, para cada prioridad que haya tenido al menos un
paciente atendido, se escribe una línea con el formato exacto
\verb|prioritat <p>: <m>|, donde \verb|<p>| es la prioridad (1, 2
o 3) y \verb|<m>| es el tiempo de espera medio con 2 decimales
(el prefijo \verb|prioritat | debe escribirse tal cual, igual que
en la versión en catalán).

\Observation

En los archivos públicos (icono del gatito) encontrarás:
\verb|main.cc| (programa principal que debes completar),
\verb|pacient.hh|, los ficheros \verb|*.hh| de los contenedores de PRO2 que puedas necesitar (incluyendo \verb|assert.hh|), un
\verb|Makefile|, \verb|.clang-format| y el directorio \verb|.vscode|.

\medskip

Debes completar el archivo \verb|main.cc| y enviarlo.

\medskip

Para poner la precisión de la salida a 2 decimales, utiliza al
principio del \verb|main|:

\begin{verbatim}
cout.setf(ios::fixed);
cout.precision(2);
\end{verbatim}

\SampleTwoCol
